%% kapa: exam=1 round=28 year=2017 period=1 form=A subject=civil_law subject_ko=민법 seq=16 number=16
%% source: Q-Net 감정평가사 기출문제 (cst00309) artlSeq=5199603 file=extracted/1008/5199603/2017년 제28회 감정평가사 1차시험 문제지/1교시 A형.pdf
%% conversion: kordoc 4.24.51 -> tools/gen1.py (auto); review=proofread-sonnet; flags=none
\begin{question}{16}
법률행위의 무효에 관한 설명으로 옳지 않은 것은? (다툼이 있으면 판례에 따름)
\begin{choices}
  \item 무효인 재산상 법률행위에 대하여 당사자가 무효임을 알고 추인하면 그 추인에는 원칙적으로 소급효가 인정된다.
  \item 위증하기로 하는 계약은 당사자가 무효임을 알고 추인하여도 유효로 될 수 없다.
  \item 불공정한 법률행위에 대하여도 무효행위의 전환에 관한 민법 규정이 적용될 수 있다.
  \item 무효행위의 추인은 명시적으로뿐만 아니라 묵시적으로도 할 수 있다.
  \item 무효인 법률행위에 따른 법률효과를 침해하는 것처럼 보이는 채무불이행이 있다고 하여도 그 법률효과의 침해에 따른 손해배상을 청구할 수는 없다.
\end{choices}
\end{question}
\answer{1}
